\section{The model and the exact law}\label{sec:model}

Fix $d\ge2$ and a generator $Q$ on $[d]=\{1,\dots,d\}$, with rates $q_{ij}\ge0$ for $i\ne j$,
exit rates $q_i=\sum_{j\ne i}q_{ij}$, $Q_{ii}=-q_i$ and $q_{\max}=\max_iq_i$. The start $\mu$ is
a probability law on $[d]$, and $\mu(F)=\sum_{i\in F}\mu_i$. For $N\ge2$ and $T>0$ put $h=T/N$,
$L=\log N$ and $\tau=T/L$. The chain $X_1,\dots,X_N$ has $X_1\sim\mu$ and transition matrix
$I+hQ$. So it switches from $i$ to $j$ with probability $hq_{ij}$ and stays with probability
$1-hq_i$. We assume $hq_{\max}\le\frac12$ unless we say otherwise, and the exact identities below
only need $hq_{\max}\le1$. The composition is $n=(n_1,\dots,n_d)$, where $n_i$ is the number of
$t\le N$ with $X_t=i$. We write $P_N(n)$ for its probability and $\alpha=n/N$. We are interested
in $T$ of order $\log N$, so that the switching probabilities are of order $(\log N)/N$.

A face is a nonempty set $F\subseteq[d]$, of size $k=k_F=\lvert F\rvert$. We use $F$ and $G$ for
faces. Let $\Delta_F$ be the set of probability vectors on $[d]$ that vanish off $F$, and
$\Delta_F^\circ$ the set of those that are positive on $F$. So $n/N\in\Delta_F^\circ$ exactly when
$\supp n=F$. For coordinates we number the states of $F$ as $1,\dots,k$ and use
$(\alpha_2,\dots,\alpha_k)$ on $\Delta_F$ and $(t_2,\dots,t_k)$ on $\{t\in\R^F:\sum_it_i=T\}$.
Any other $k-1$ coordinates give the same Lebesgue measure. For $x,y\in\R^F$ we write
$\langle x,y\rangle=\sum_{i\in F}x_iy_i$ and $x\circ y=(x_iy_i)_{i\in F}$. Let
$A_F=(q_{ij})_{i,j\in F}$ with zero diagonal, and let $Q_F=A_F-\operatorname{diag}(q_i)_{i\in F}$
be the principal submatrix of $Q$ on $F$. Note that $q_i$ includes the rates of leaving $F$, so
$Q_F$ generates the chain killed when it leaves $F$. The switching graph $G_F$ is the directed
graph on $F$ with an edge $i\to j$ whenever $i\ne j$ and $q_{ij}>0$. A matrix is Metzler if its
off-diagonal entries are nonnegative. The Perron root $\varrho(B)$ of a Metzler matrix $B$ is its
largest real eigenvalue. It exists, and if $B$ is irreducible it is a simple eigenvalue with
positive left and right eigenvectors, by the Perron--Frobenius theorem for $B+cI\ge0$ with $c$
large~\cite{bermanplemmons1994}.

\begin{definition}[faces]\label{def:faces}
Let $F$ be a face. We call $F$ \emph{irreducible} if $G_F$ is strongly connected, so every
singleton is irreducible. We call $F$ \emph{admissible} if for some $N$ some path
$x_1\cdots x_N$ of positive probability has support exactly $F$. We say that $F$ satisfies
\textup{(H)} if $F$ is irreducible and $\mu(F)>0$. The \emph{exit rate} of $F$ is
$\lambda_F=-\varrho(Q_F)$, so $\lambda_{\{i\}}=q_i$. If $F$ is irreducible, $r_F$ and $l_F$ are
right and left eigenvectors of $Q_F$ for $-\lambda_F$ with positive entries and
$\langle l_F,r_F\rangle=1$, and $\alpha^*_F=l_F\circ r_F$ is the \emph{quasi-ergodic law} of
$F$~\cite{darrochseneta1965,darrochseneta1967}. The \emph{face maximum} is
$M_F=\max\{P_N(n):\supp n=F\}$.
\end{definition}

The pair $(r_F,l_F)$ is unique up to $(cr_F,l_F/c)$, and $\alpha^*_F\in\Delta_F^\circ$ does not
depend on $c$. When $hq_{\max}<1$, a path has positive probability exactly when $\mu_{x_1}>0$ and
$q_{x_tx_{t+1}}>0$ at every switch, and repeating its last state keeps it positive. So $M_F>0$ for
all large $N$ if $F$ is admissible, and $M_F=0$ otherwise. If $F$ is reducible, we order the
strongly connected components $C_1,\dots,C_r$ of $G_F$ so that every edge between different
components goes forward. Then $Q_F$ is block triangular with diagonal blocks $Q_{C_j}$ (its
Frobenius normal form~\cite{bermanplemmons1994}), so $\lambda_F=\min_j\lambda_{C_j}$. If $F$ is
irreducible, the row sums of $e^{TQ_F}\ge0$ are the probabilities that the continuous-time chain
started in $F$ stays in $F$ up to time $T$, and $e^{TQ_F}r_F=e^{-T\lambda_F}r_F$. Comparing $\ones$
with $r_F$ shows that these probabilities are $e^{-T\lambda_F}$ up to a factor between
$\min r_F/\max r_F$ and its inverse. So $\lambda_F$ is the rate at which the chain leaves $F$.

\begin{lemma}[admissible faces]\label{lem:admissible}
Let $F$ be a face and $C_1,\dots,C_r$ the strongly connected components of $G_F$. \textup{(a)} $F$
is admissible if and only if the components can be numbered so that $\mu(C_1)>0$ and, for each $j<r$,
$q_{xy}>0$ for some $x\in C_j$ and $y\in C_{j+1}$. \textup{(b)} An irreducible face is admissible
if and only if $\mu(F)>0$. So \textup{(H)} holds exactly for the irreducible admissible faces.
\textup{(c)} If $Q$ has symmetric support, that is $q_{ij}>0$ exactly when $q_{ji}>0$, then every
admissible face is irreducible.
\end{lemma}

\begin{proof}
The states of the maximal runs of a path of positive probability with support $F$ form a walk in
$G_F$ that starts in $\supp\mu$ and visits every state of $F$. For $N$ at least its length, every
such walk arises in this way. So $F$ is admissible if and only if $G_F$ has such a walk. (a) Along
such a walk, a component $C$ cannot be visited again after a state $y\notin C$, since then $y$ and
$C$ would be reachable from each other. So the walk passes through each component in one block.
The first block contains the start, and consecutive blocks are joined by an edge.
Conversely, given the numbering and edges $x_j\to y_{j+1}$ from $C_j$ to $C_{j+1}$, take
$y_1\in C_1$ with $\mu_{y_1}>0$ and any $x_r\in C_r$. Since $C_j$ is strongly connected, some walk
from $y_j$ to $x_j$ visits every state of $C_j$. Joining these walks by the edges
$x_j\to y_{j+1}$ gives the required walk. (b) is (a) with $r=1$. (c) With symmetric support an
edge from $C_1$ to $C_2$ gives an edge back, which would merge $C_1$ and $C_2$. So (a) forces
$r=1$.
\end{proof}

For a face $F$ and $m\in\Z_{\ge1}^F$ put $M=\sum_{i\in F}m_i$ and
$W(m)=\sum_c\mu_{c_1}\prod_{j<M}q_{c_jc_{j+1}}$, where the sum runs over the words
$c=c_1\cdots c_M$ in $F$ with $c_j\ne c_{j+1}$ for $j<M$ in which each $i\in F$ occurs $m_i$ times.
We do not show the dependence on $F$ in $W$, or in $H$ and $f_T$ below.

\begin{lemma}[run expansion]\label{lem:runs}
Let $hq_{\max}\le1$ and let $n$ have support $F$. Then, with $\binom{n_i-1}{m_i-1}=0$ for
$m_i>n_i$ and $0^0=1$,
\[
P_N(n)=\sum_{m\in\Z_{\ge1}^F}W(m)\,h^{M-1}\prod_{i\in F}\binom{n_i-1}{m_i-1}(1-hq_i)^{n_i-m_i}.
\]
\end{lemma}

\begin{proof}
Split a path with support $F$ into maximal runs. Their states $c_1\cdots c_M$ form a word as in
$W(m)$, where $m_i$ is the number of runs of state $i$. The path has $M-1$ switches, of
probabilities $hq_{c_jc_{j+1}}$, and $n_i-m_i$ stays in state $i$, of probability $1-hq_i$ each.
So its probability is $\mu_{c_1}h^{M-1}\prod_{j<M}q_{c_jc_{j+1}}\prod_{i\in F}(1-hq_i)^{n_i-m_i}$.
For a fixed $c$ the path is determined by the run lengths, and the lengths of the runs of state
$i$ form one of the $\binom{n_i-1}{m_i-1}$ compositions of $n_i$ into $m_i$ positive parts,
independently over $i$. Summing over $c$ gives the claim.
\end{proof}

The continuous-time analogue replaces the binomial coefficients by powers. For
$t\in(0,\infty)^F$ put
$H(t)=\sum_{m\in\Z_{\ge1}^F}W(m)\prod_{i\in F}t_i^{m_i-1}/(m_i-1)!$. There are at most $k^M$
words of length $M$, so $W(m)\le k^Mq_{\max}^{M-1}$ and
$H(t)\le k^kq_{\max}^{k-1}e^{kq_{\max}\sum_it_i}<\infty$.

\begin{lemma}[the continuum density]\label{lem:continuum}
Let $(Y_s)_{s\ge0}$ be the continuous-time chain with generator $Q$ and $Y_0\sim\mu$, let
$\mathcal O_i=\int_0^T\mathbf 1\{Y_s=i\}\,ds$, and let $E_F$ be the event that $Y$ stays in $F$
during $[0,T]$ and visits every state of $F$. Then $\Prob(E_F,\ \mathcal O\in\cdot\,)$ has the
density $e^{-\langle q,t\rangle}H(t)$ in $(t_2,\dots,t_k)$ on $\{t\in(0,\infty)^F:\sum_it_i=T\}$,
and $\Prob(E_F,\ \mathcal O/T\in\cdot\,)$ has the density
\[
f_T(\alpha)=T^{k-1}e^{-T\langle q,\alpha\rangle}H(T\alpha)
\]
in $(\alpha_2,\dots,\alpha_k)$. For $n$ with support $F$ and $t=hn$,
$N^{-(k-1)}f_T(n/N)=h^{k-1}e^{-\langle q,t\rangle}H(t)$. For $F=\{i\}$,
$f_T=\Prob(E_F)=\mu_ie^{-q_iT}$.
\end{lemma}

\begin{proof}
For a word $c=c_1\cdots c_M$ in $F$ with $c_j\ne c_{j+1}$ that contains every state of $F$, let
$E_c$ be the event that $Y$ visits exactly $c_1,\dots,c_M$ in this order during $[0,T]$. Since $Y$
jumps finitely often in $[0,T]$ almost surely, $E_F$ is the disjoint union of the events $E_c$ up
to a null set. Fix $c$, let $s_1,\dots,s_{M-1}$ be the first $M-1$ holding times and
$s_M=T-\sum_{j<M}s_j$. The holding time in $c_j$ is exponential with rate $q_{c_j}$, the jump to
$c_{j+1}$ has probability $q_{c_jc_{j+1}}/q_{c_j}$, and the last holding time is censored. It
only has to exceed $s_M$, which has probability $e^{-q_{c_M}s_M}$. So on
$\{s_j>0\text{ for }j\le M\}$ the event $E_c$ has density
$\mu_{c_1}\prod_{j<M}q_{c_jc_{j+1}}e^{-\langle q,\mathcal O\rangle}$ in $(s_1,\dots,s_{M-1})$,
where $\mathcal O_i=\sum_{j:c_j=i}s_j$. (If $q_{c_j}=0$ for some $j<M$, then $E_c$ is null and the formula gives $0$.) Let
$j_i$ be the position of the last letter $i$ in $c$, so $j_{c_M}=M$. Replacing $s_{j_i}$ by
$\mathcal O_i$ for $i\ne c_M$ is a triangular linear change of variables with unit diagonal. For
fixed $\mathcal O=t$ the other $m_i-1$ holding times in state $i$ range over
$\{x\in(0,\infty)^{m_i-1}:\sum x<t_i\}$, for every $i\in F$. For $i=c_M$ this is the condition
$s_M>0$, with $t_{c_M}=T-\sum_{i\ne c_M}t_i$. The set has volume $t_i^{m_i-1}/(m_i-1)!$ and is
empty unless $t_i>0$. Integrating out these variables gives the density of $\mathcal O$ on $E_c$
in $(t_i)_{i\ne c_M}$, and hence in $(t_2,\dots,t_k)$. Summing over $c$ by monotone convergence
gives $e^{-\langle q,t\rangle}H(t)$. The rest follows from the change of variables $t=T\alpha$, with Jacobian
$T^{k-1}$, and from $T/N=h$. For $k=1$ the only word is $c=i$, so $H=\mu_i$.
\end{proof}

\begin{example}[the 4 models]\label{ex:models}
(i) For $d=2$, $q_{12}=q_{21}=1$ and uniform $\mu$ the chain is the persistent walk
of~\cite{paperA} with $p=1-h$, so $T=N(1-p)$. Bin $k$ is the composition $(k,N-k)$, and
Lemma~\ref{lem:runs} is the run count of~\cite[Thm.~2.1]{paperA}. (ii) For $Q=J-dI$, with $J$
the all-ones matrix, and uniform $\mu$ the chain is the model of~\cite{paperB} with its
parameters $r=h$ and $p=1-(d-1)h$. So $T=Nr$, and its crossing variable $Nu=Nr/p$ is
$T(1+O(T/N))$. Here $\lambda_F=d-\lvert F\rvert$ and $\alpha^*_F$ is uniform on $F$. (iii) For
$Q=\Pi-I$, with $\Pi=(\pi_{ij})$ stochastic with zero diagonal, $I+hQ=(1-h)I+h\Pi$ and
$\lambda_F=1-\varrho(\Pi_F)$, where $\Pi_F$ is the principal submatrix of $\Pi$ on $F$. This is
the sticky chain of Section~\ref{sec:sticky}. (iv) Let $d=4$, number the directions east, north,
west and south as $0,1,2,3$ modulo $4$, and put $q_{j,j\pm1}=r$, $q_{j,j+2}=s$ and $\mu$ uniform.
Then $X_t$ is the direction of step $t$ of a walk on $\Z^2$ that turns left or right with
probability $rh$ each and reverses with probability $sh$. Its position $(n_0-n_2,n_1-n_3)$ is a
linear function of the composition (Section~\ref{sec:planar}). In (i), (ii) and (iv), $Q$ has
symmetric support and $\mu$ has full support. So by Lemma~\ref{lem:admissible} the admissible
faces are the irreducible ones, and they all satisfy \textup{(H)}.
\end{example}

\begin{remark}[the power of $N$]\label{rem:lattice-term}
Since $\sum_{i\in F}(m_i-1)=M-k$, we can write $h^{M-1}=h^{k-1}\prod_ih^{m_i-1}$ in
Lemma~\ref{lem:runs}. With $t=hn$,
$h^{m_i-1}\binom{n_i-1}{m_i-1}=\frac{t_i^{m_i-1}}{(m_i-1)!}\prod_{l<m_i}\bigl(1-\frac l{n_i}\bigr)$,
and $(1-hq_i)^{n_i-m_i}$ is close to $e^{-q_it_i}$. So each term with every $m_i^2$ small compared
with $n_i$ is close to the matching term of
$h^{k-1}e^{-\langle q,t\rangle}H(t)=N^{-(k-1)}f_T(n/N)$. The terms that matter have $M$ of order
$T$, and under \textup{(H)} Lemma~\ref{lem:discrete} shows that the relative error is $O(T^2/N)$
on compact subsets of $\Delta_F^\circ$. So a run after the first in state $i$ brings a factor $h$
for its switch and a factor of order $n_i$ for its placement, and together these give a factor of
order $T$. Only the $k-1$ switches that enter a new state leave a factor $h=T/N$ each.
\end{remark}
